\documentclass[aspectratio=169]{beamer}

% ---------- ENCODING & FONTS ----------
\usepackage[utf8]{inputenc}
\usepackage[T1]{fontenc}

% ---------- THEME & LOOK ----------
\usetheme{Madrid}
\usecolortheme{seahorse}
\usefonttheme{professionalfonts}
\setbeamertemplate{navigation symbols}{}
\setbeamercovered{transparent}
\setbeamertemplate{itemize items}[circle]
\setbeamertemplate{enumerate items}[default]
\setlength{\parskip}{0.25em}

% ---------- COLORS ----------
\usepackage{xcolor}
\definecolor{accent}{RGB}{0,92,184}
\definecolor{soft}{RGB}{233,244,255}
\setbeamercolor{structure}{fg=accent}
\setbeamercolor{title}{fg=white,bg=accent}
\setbeamercolor{frametitle}{fg=accent}
\setbeamercolor{block title}{bg=accent!12,fg=accent}
\setbeamercolor{block body}{bg=soft}

% ---------- PACKAGES ----------
\usepackage{booktabs}
\usepackage{amsmath, amssymb}
\usepackage{array}
\usepackage{multirow}
\usepackage{hyperref}
\hypersetup{colorlinks=true,linkcolor=accent,urlcolor=accent}

% ---------- SAFE FALLBACK FOR TRANSITIONS ----------
\makeatletter
\@ifundefined{transfade}{\newcommand{\transfade}[1][]{}}{}
\makeatother

% ---------- PLACEHOLDER BOX (compile-safe; no images required) ----------
\newcommand{\placeholderbox}[2][0.82\linewidth]{%
  \begin{center}
    \fbox{\begin{minipage}[c][0.46\textheight][c]{#1}
      \centering \vspace{0.25em}
      \textbf{#2}\\[0.6em]
      \emph{Insert figure/diagram here later.}
    \end{minipage}}
  \end{center}
}

% ---------- TITLE ----------
\title{\textbf{Labour Supply over the Life Cycle \vspace{.5cm}\\  }}
\subtitle{ {\large ECON 300 - Labour Economics}}
%\institute{UNBC}
\author{Muhebullah Karimzada}

\date{Winter 2026}



\begin{document}

% ============================================================
% TITLE
% ============================================================
\begin{frame}[plain]
  \titlepage
\end{frame}

% ============================================================
\section{Learning Objectives \& Main Questions}
% ============================================================

\begin{frame}{Learning Objectives}
\transfade
\begin{itemize}[<+->]
  \item Explain why cross-sectional age comparisons can be misleading.
  \item Distinguish the effects of \textbf{temporary} vs.\ \textbf{permanent} wage changes on labour supply.
  \item Describe how economists model household choices over paid work, unpaid work, and fertility.
  \item Outline Canada’s three-tier pension system and retirement incentives.
  \item Use diagrams to show how public pensions affect retirement timing and income.
\end{itemize}
\end{frame}

\begin{frame}{Main Questions}
\transfade
\begin{itemize}[<+->]
  \item Do men’s and women’s labour supply patterns evolve similarly with age?
  \item How do life-cycle planning horizons change labour-supply responses to wages?
  \item What new complexities arise when modelling each period of life?
  \item How do labour-saving household technologies reshape time allocation?
  \item How do fertility choices interact with women’s labour supply?
  \item What determines retirement age, and how do pensions shape it?
\end{itemize}
\end{frame}

% ============================================================
\section{Characterizing Life-Cycle Labour Supply}
% ============================================================

\begin{frame}{Age--Participation Profiles}
\transfade
\begin{itemize}[<+->]
  \item Participation rates vary systematically with age for both men and women.
  \item Women’s age profiles have shifted upward over successive cohorts.
  \item Interpretation requires disentangling \alert{age}, \alert{period}, and \alert{cohort} effects.
\end{itemize}
\end{frame}

\begin{frame}{Figure 4.1 — Age--Labour Force Participation Profiles \textit{(Placeholder)}}
\transfade
\placeholderbox{Figure 4.1 — Plot age-participation profiles for men and women; show cohort shifts.}
\end{frame}

\begin{frame}{Cohort Effect: Definition}
\transfade
\begin{itemize}[<+->]
  \item The difficulty of separating the influence of \textbf{age} from \textbf{year of birth} (vintage) at a point in time.
  \item For women, newer cohorts tend to be more attached to the labour market than earlier cohorts.
\end{itemize}
\end{frame}

\begin{frame}{Worked Example — What Is a Cohort Effect?}
\transfade
\small
\begin{tabular}{lccc}
\toprule
 & \textbf{Observed in 2000} & \textbf{2005} & \textbf{2010} \\
\midrule
\textbf{Older Cohort, born in 1975:} & & & \\
\quad Age & 25 & 30 & 35 \\
\quad Participation Rate & 75\% & 80\% & 85\% \\
\addlinespace[0.25em]
\textbf{Younger Cohort, born 1980:} & & & \\
\quad Age & 20 & 25 & 30 \\
\quad Participation Rate & 75\% & 80\% & 85\% \\
\bottomrule
\end{tabular}

\vspace{0.4em}
\footnotesize
\textbf{Takeaway:} Cross-section comparisons can confound true aging patterns with cohort differences. Tracking cohorts over time isolates life-cycle effects.
\end{frame}

% ============================================================
\section{Dynamic Life-Cycle Model}
% ============================================================

\begin{frame}{Two Modelling Approaches}
\transfade
\begin{itemize}[<+->]
  \item \textbf{Sequence of static problems:} convenient but ignores forward-looking behaviour.
  \item \textbf{Dynamic framework:} individuals choose consumption and leisure for each of \(N\) periods to maximize lifetime utility given an intertemporal budget constraint.
\end{itemize}
\end{frame}

\begin{frame}{Assumptions \& Setup}
\transfade
\begin{itemize}[<+->]
  \item Horizon of \(N\) periods; preferences over \((c_t,\ell_t)_{t=1}^N\).
  \item For intuition, start with no nonlabour income and perfect foresight.
  \item Labour supply each period links through \textbf{present-value} budget constraint.
\end{itemize}
\end{frame}

\begin{frame}{Lifetime Income Constraint (Two-Period Illustration)}
\transfade
\small
Income each period: \(y_t = w_t h_t\), with \(h_t = T - \ell_t\).\par
Let real interest rate be \(r\). Present-value constraint:
\[
\text{PV}_0(c) \;=\; \sum_{t=1}^{2}\frac{c_t}{(1+r)^{t-1}}
\;=\; \sum_{t=1}^{2}\frac{w_t (T-\ell_t)}{(1+r)^{t-1}} \;=\; \text{PV}_0(y).
\]
Equivalently, replacing \(h_t\) by \(\ell_t\): higher wages raise the opportunity cost of leisure across periods.
\end{frame}

\begin{frame}{Figure 4.2 — Dynamic Life-Cycle Wage Changes \textit{(Placeholder)}}
\transfade
\placeholderbox{Figure 4.2 — Contrast temporary vs.\ permanent and anticipated vs.\ unanticipated wage changes; show income vs.\ substitution effects.}
\end{frame}

\begin{frame}{Temporary vs.\ Permanent Wage Increases}
\transfade
\begin{itemize}[<+->]
  \item \textbf{Temporary} increase: strong substitution effect in that period, modest income effect.
  \item \textbf{Permanent} increase: larger lifetime income effect; labour supply response distributed across periods.
  \item Anticipation matters: planned time shifts differ if change is known in advance.
\end{itemize}
\end{frame}

% ============================================================
\section{Household Production}
% ============================================================

\begin{frame}{Households as Producers and Consumers}
\transfade
\begin{itemize}[<+->]
  \item Households allocate scarce time among \textbf{pure leisure}, \textbf{household work}, and other activities.
  \item Utility from consumption often requires time inputs (e.g., cooking).
\end{itemize}
\end{frame}

\begin{frame}{Basic Framework}
\transfade
\small
Utility over produced goods \(Z_i\):\quad \(U=f(Z_1,\dots,Z_n)\), \quad
\(Z_i=g_i(X_1,\dots,X_m; h_i)\)\par
\begin{itemize}[<+->]
  \item \(X\): market-purchased inputs; \(h_i\): time spent producing \(Z_i\).
  \item \textbf{Leisure} is not purely separate; time is an input into many \(Z_i\).
\end{itemize}
\end{frame}

\begin{frame}{Worked Example — Food at Home vs.\ Precooked}
\transfade
\small
Utility \(u=f(\text{nutrition}, \text{rest})\).\par
\begin{itemize}[<+->]
  \item Prices affecting choices:
  \begin{itemize}[<+->]
    \item Price of leisure \(= W\) (wage).
    \item Unit nutrition from raw food: \(P_R + Wd\) (goods price + time).
    \item Unit nutrition from precooked food: \(P_C\).
  \end{itemize}
  \item \textbf{Comparative statics:}
  \begin{itemize}[<+->]
    \item \(\downarrow P_C\): more precooked, more leisure; effect on raw food ambiguous.
    \item \(\downarrow P_R\): more raw food, less leisure; effect on precooked ambiguous.
    \item \(\uparrow W\): tilt away from time-intensive goods.
    \item \(\downarrow d\): cheaper home production; leisure/labour supply depends on income vs.\ substitution.
  \end{itemize}
\end{itemize}
\end{frame}

% ============================================================
\section{Fertility and Childbearing}
% ============================================================

\begin{frame}{Determinants of Fertility}
\transfade
\begin{itemize}[<+->]
  \item \textbf{Income:} higher income can raise desired children (income effect), but higher potential earnings raise the opportunity cost of time (substitution effect).
  \item \textbf{Cost of children:} foregone earnings, direct expenses.
  \item \textbf{Prices of related goods:} daycare, education, medical care.
  \item \textbf{Tastes and preferences:} social norms, family planning, cultural and religious factors.
  \item \textbf{Technology:} contraceptives and labour-saving home production.
\end{itemize}
\end{frame}

% ============================================================
\section{Retirement Decisions and Pensions}
% ============================================================

\begin{frame}{Retirement: Definitions and Policy Relevance}
\transfade
\begin{itemize}[<+->]
  \item Retirement can mean leaving the labour force, reducing hours, or moving to less demanding work.
  \item Implications for private savings, unemployment, labour force size, and national income.
\end{itemize}
\end{frame}

\begin{frame}{Theoretical Determinants of Retirement}
\transfade
\begin{itemize}[<+->]
  \item Mandatory retirement provisions (compulsory vs.\ automatic).
  \item Wealth and earnings:
    \begin{itemize}[<+->]
      \item \(\uparrow\) Wealth \(\Rightarrow\) income effect \(\Rightarrow\) more leisure (earlier retirement).
      \item \(\uparrow\) Earnings: income and substitution effects offset; net effect ambiguous.
    \end{itemize}
  \item Health, job difficulty, and family composition (e.g., dual earners).
\end{itemize}
\end{frame}

\begin{frame}{Canada’s Pension System: Three Tiers}
\transfade
\begin{itemize}[<+->]
  \item \textbf{Old Age Security (OAS):} tax-financed demogrant at 65; GIS supplements are means-tested.
  \item \textbf{CPP/QPP:} social insurance funded by compulsory employer/employee contributions; benefits linked to prior earnings.
  \item \textbf{Employer Pensions and Private Savings:} defined benefit/contribution plans, RRSPs, and other private arrangements.
\end{itemize}
\end{frame}

\begin{frame}{Formal Income--Leisure Model with Pension}
\transfade
\small
\[
Y_p = B + (1 - p - t)\,W\,(T - \ell),
\]
where \(Y_p\): income of a pension receiver; \(B\): pension; \(W\): wage; \(T\): time endowment; \(\ell\): leisure; \(p\): payroll tax rate; \(t\): implicit tax (clawback).
\end{frame}

\begin{frame}{Figure 4.3 — Budget Constraints under Social Insurance Pensions \textit{(Placeholder)}}
\transfade
\placeholderbox{Figure 4.3 — Budget lines with pension \(B\) and implicit tax/clawback; assume payroll tax \(p=0\) for illustration.}
\end{frame}

\begin{frame}{Employer-Sponsored Pensions: Types}
\transfade
\begin{itemize}[<+->]
  \item Flat benefit, earnings-based benefit, and defined-contribution plans.
  \item \textbf{Other provisions:} backloading; early/special retirement; postponed retirement penalties.
\end{itemize}
\end{frame}

\begin{frame}{Backloading and Retirement Incentives}
\transfade
\begin{itemize}[<+->]
  \item Benefits increase more steeply with seniority; encourages retention.
  \item Young workers: incentive to stay; older workers: incentive to delay retirement \emph{until} vesting thresholds, then retire.
\end{itemize}
\end{frame}

\begin{frame}{Early/Special vs.\ Postponed Retirement}
\transfade
\begin{itemize}[<+->]
  \item Early retirement: often at age \(\sim55\) with 10+ years of service (less generous due to shorter service).
  \item Special provisions: ages \(60\)–\(62\) with 20+ years; common in earnings-based plans; extensively subsidized.
  \item Postponed retirement: penalties beyond normal age in some plans.
\end{itemize}
\end{frame}

\begin{frame}{Picking the Perfect Retirement Age \textit{(Placeholder)}}
\transfade
\placeholderbox{Decision Map — Weigh pension wealth, health, job strain, spouse’s work, and implicit taxes.}
\end{frame}

% ============================================================
\section{Chapter Summary}
% ============================================================

\begin{frame}{Summary}
\transfade
\begin{itemize}[<+->]
  \item Life-cycle labour supply shows systematic age patterns; women’s profiles have shifted up across cohorts.
  \item Dynamic models link period-by-period labour supply through a present-value constraint.
  \item Household production embeds time costs into consumption; technology and wages tilt choices.
  \item Fertility responds to income, prices of related goods, tastes, and technology, interacting with women’s labour supply.
  \item Canada’s pensions (OAS, CPP/QPP, employer/private plans) and plan design shape retirement timing and incomes.
\end{itemize}
\end{frame}

\begin{frame}
\centering
\Huge{\textbf{Thank You!}} \\

\end{frame}

\end{document}
